\documentclass[9pt,landscape,a4paper]{article}
\usepackage[utf8]{inputenc}
\usepackage[T1]{fontenc}
\usepackage{geometry}
\usepackage{multicol}
\usepackage{amsmath,amssymb}
\usepackage{enumitem}
\usepackage{graphicx}
\usepackage{xcolor}
\usepackage{listings}
\usepackage{booktabs}
\usepackage{titlesec}
\usepackage{fancyhdr}
\usepackage{hyperref}
\usepackage{CJKutf8}

% Page layout
\geometry{top=0.7cm, bottom=0.7cm, left=0.7cm, right=0.7cm}
\pagestyle{empty}
\setlength{\parindent}{0pt}
\setlength{\parskip}{1.5pt}
\setlength{\columnsep}{7pt}

% Section formatting
\titleformat{\section}{\normalsize\bfseries\color{blue!70!black}}{}{0em}{}[\titlerule]
\titleformat{\subsection}{\small\bfseries\color{blue!50!black}}{}{0em}{}
\titlespacing{\section}{0pt}{4pt}{2pt}
\titlespacing{\subsection}{0pt}{3pt}{1pt}

% Code listing style
\lstset{
    language=Python,
    basicstyle=\ttfamily\fontsize{6.5}{7.5}\selectfont,
    keywordstyle=\color{blue!70!black}\bfseries,
    stringstyle=\color{red!60!black},
    commentstyle=\color{green!50!black}\itshape,
    numberstyle=\tiny\color{gray},
    backgroundcolor=\color{gray!8},
    frame=single,
    framerule=0.3pt,
    rulecolor=\color{gray!50},
    breaklines=true,
    breakatwhitespace=false,
    tabsize=2,
    showstringspaces=false,
    aboveskip=2pt,
    belowskip=2pt,
    xleftmargin=2pt,
    xrightmargin=2pt,
}

\newcommand{\keybox}[1]{%
    \begin{center}
    \fcolorbox{blue!50!black}{blue!5}{%
        \parbox{0.95\linewidth}{\small #1}%
    }
    \end{center}
}

\newcommand{\zhint}[1]{{\scriptsize\color{gray!70!black}#1}}

\begin{document}
\begin{CJK}{UTF8}{gbsn}

\begin{center}
{\large\bfseries\color{blue!70!black} Chapter 8 -- Data Visualization}\\
{\footnotesize IERG4320/IEMS5726 Data Science in Practice | Cheat Sheet}
\end{center}

\begin{multicols}{3}

% ==================== SECTION 1 ====================
\section{Introduction}
\zhint{数据可视化：用图形展示数据中的模式和趋势}

\textbf{Goals}: explore, monitor, explain data.\\
\textbf{Benefits}: reveal patterns/trends, identify outliers, simplify complex info, aid decision-making.

\textbf{When to visualize}:
\begin{itemize}[leftmargin=8pt,itemsep=0pt,topsep=0pt]
\item Preprocessing: distribution, trends \zhint{看数据分布}
\item Analysis: model performance \zhint{看模型表现}
\item Interpretation: present findings \zhint{展示结果}
\end{itemize}

\textbf{Packages}: \texttt{matplotlib.pyplot as plt}, \texttt{seaborn as sns}\\
\zhint{Seaborn建在matplotlib之上，和pandas配合好}

% ==================== SECTION 2 ====================
\section{Preprocessing Visualizations}

\subsection{1. Histogram (Distribution)}
\zhint{直方图：看单个特征的数据分布}
\begin{lstlisting}
import seaborn as sns
import matplotlib.pyplot as plt
penguins = sns.load_dataset("penguins")
# Basic histogram
sns.displot(penguins, x="flipper_length_mm")
# Custom bins
sns.displot(penguins, x="flipper_length_mm",
            binwidth=3)
sns.displot(penguins, x="flipper_length_mm",
            bins=20)
# Colored by category
sns.displot(penguins, x="flipper_length_mm",
            hue="species")
# Stacked
sns.displot(penguins, x="flipper_length_mm",
  hue="species", multiple="stack")
# Side-by-side
sns.displot(penguins, x="flipper_length_mm",
  hue="sex", multiple="dodge")
plt.show()
\end{lstlisting}

\subsection{2. Boxplot}
\zhint{箱线图：显示中位数、四分位数和异常值}\\
\textbf{Key}: Median, Q1 (25\%), Q3 (75\%), IQR=Q3--Q1.\\
Whiskers: Q1--1.5$\times$IQR to Q3+1.5$\times$IQR. Beyond = outliers.

\begin{lstlisting}
# Pandas approach
penguins.boxplot(column=["flipper_length_mm"])
# Conditional boxplot
penguins.boxplot(column=["flipper_length_mm"],
                 by=["species"])
# Seaborn approach
sns.boxplot(data=penguins,
  x="flipper_length_mm", y="sex",
  hue="species")
plt.show()
\end{lstlisting}

\subsection{3. Scatter Plot}
\zhint{散点图：看两个特征的关系}
\begin{lstlisting}
# Matplotlib
plt.scatter(x, y)
plt.title("Title"); plt.xlabel("X")
plt.ylabel("Y"); plt.savefig("name")
# Seaborn (with hue grouping)
sns.scatterplot(x="total_bill", y="tip",
  data=tips, hue="smoker",
  hue_order=["Yes","No"])
plt.show()
\end{lstlisting}

\subsection{4. Q-Q Plot}
\zhint{QQ图：检查数据是否服从某分布（如正态分布）}\\
Points on 45-degree line = same distribution.
\begin{lstlisting}
import statsmodels.api as sm
import pylab as py
data_pts = np.random.normal(0, 1, 100)
sm.qqplot(data_pts, line='45')
py.show()
\end{lstlisting}

\subsection{5. Heatmap}
\zhint{热力图：用颜色表示矩阵中数值大小，常用于看特征相关性}
\begin{lstlisting}
glue = sns.load_dataset("glue").pivot(
  index="Model",columns="Task",values="Score")
sns.heatmap(glue)           # basic
sns.heatmap(glue, annot=True) # with numbers
plt.show()
\end{lstlisting}

\subsection{6. Pair Plot}
\zhint{配对图：所有特征两两组合的散点图+对角线分布图}
\begin{lstlisting}
sns.pairplot(USAhousing)
plt.show()
\end{lstlisting}

% ==================== SECTION 3 ====================
\section{Analysis Visualizations}

\subsection{1. ROC Curve}
\zhint{ROC曲线：FPR vs TPR，AUC越大模型越好}
\begin{lstlisting}
from sklearn.metrics import roc_curve, \
  roc_auc_score
fpr, tpr, _ = roc_curve(y_test, y_proba)
auc = roc_auc_score(y_test, y_proba)
plt.plot(fpr, tpr, label=f"AUC={auc:.3f}")
plt.xlabel('False Positive Rate')
plt.ylabel('True Positive Rate')
plt.legend(); plt.show()
\end{lstlisting}

\subsection{2. Training Curve (ANN)}
\zhint{训练曲线：看loss是否收敛，train和val差距是否合理}
\begin{lstlisting}
plt.plot(iters, losses, label="Train")
plt.xlabel("Iterations")
plt.ylabel("Loss"); plt.show()
# Accuracy curve
plt.plot(iters, train_acc, label="Train")
plt.plot(iters, val_acc, label="Validation")
plt.legend(loc='best'); plt.show()
\end{lstlisting}
\zhint{train和val差距大=过拟合; 都高=欠拟合}

\subsection{3. Predicted vs Actual}
\zhint{预测值vs真实值：点越接近对角线=模型越好}
\begin{lstlisting}
plt.scatter(true_value, predicted_value,
            c='crimson')
p1=max(max(predicted_value),max(true_value))
p2=min(min(predicted_value),min(true_value))
plt.plot([p1,p2],[p1,p2],'b-') # diagonal
plt.xlabel('True'); plt.ylabel('Pred')
plt.show()
\end{lstlisting}

\subsection{4. Residual Plot}
\zhint{残差图：检查回归模型假设是否成立（残差应随机分布）}
\begin{lstlisting}
sns.residplot(x='Head_size',
  y='Brain_weight', data=data)
plt.show()
\end{lstlisting}

\subsection{5. Calibration Plot}
\zhint{校准图：预测概率是否和实际频率一致}
\begin{lstlisting}
from sklearn.calibration import calibration_curve
prob_true, prob_pred = calibration_curve(
  y_test, y_probs, n_bins=10)
plt.plot(prob_pred, prob_true, marker='o')
plt.plot([0,1],[0,1],'--',color='gray')
plt.xlabel('Mean predicted probability')
plt.ylabel('Fraction of positives')
plt.show()
\end{lstlisting}

\subsection{6. Learning Curve}
\zhint{学习曲线：看增加训练数据是否能提升模型}
\begin{lstlisting}
from sklearn.model_selection import learning_curve
sizes, train_sc, val_sc = learning_curve(
  estimator=LinearRegression(),
  X=df[features], y=df[target],
  train_sizes=[1,100,500,2000,5000,7654],
  cv=5, scoring='neg_mean_squared_error')
plt.plot(sizes,-train_sc.mean(axis=1),
         label='Train error')
plt.plot(sizes,-val_sc.mean(axis=1),
         label='Validation error')
plt.legend(); plt.show()
\end{lstlisting}
\zhint{两条线差距很小=理想; 差距大=模型太简单或数据不够}

% ==================== SECTION 4 ====================
\section{Combining Multiple Plots}

\subsection{1. Subplots Grid}
\begin{lstlisting}
fig, axes = plt.subplots(nrows=2, ncols=2,
                         figsize=(12, 8))
axes[0,0].plot(x, y1); axes[0,0].set_title('A')
axes[0,1].plot(x, y2); axes[1,0].plot(x, y3)
plt.tight_layout(); plt.show()
\end{lstlisting}

\subsection{2. Overlay on Same Axes}
\begin{lstlisting}
plt.plot(x, y1, 'b-', label='sin(x)')
plt.plot(x, y2, 'r--', label='cos(x)')
plt.legend(); plt.show()
\end{lstlisting}

\subsection{3. Mixed Plot Types}
\begin{lstlisting}
fig, ((ax1,ax2),(ax3,ax4)) = plt.subplots(2,2)
ax1.bar(cats, vals)     # bar
ax2.scatter(sx, sy)     # scatter
ax3.hist(data, bins=15) # histogram
ax4.pie(vals, labels=cats, autopct='%1.1f%%')
\end{lstlisting}

\subsection{4. Shared Axes}
\begin{lstlisting}
fig, (ax1,ax2,ax3) = plt.subplots(1, 3,
  figsize=(15,5), sharey=True)
\end{lstlisting}
\zhint{sharey=True让多个子图共享Y轴，便于比较}

% ==================== SECTION 5 ====================
\section{Text Visualization}

\subsection{Word Cloud}
\zhint{词云：词频越高，字体越大}
\begin{lstlisting}
from wordcloud import WordCloud
# From text
wc = WordCloud(width=800, height=400,
  background_color='white').generate(text)
plt.imshow(wc, interpolation='bilinear')
plt.axis('off'); plt.show()

# From frequency dict
word_freq = {'python':45, 'data':38, ...}
wc = WordCloud(width=800, height=400,
  background_color='lightblue', colormap='cool'
).generate_from_frequencies(word_freq)

# Multiple documents as subplots
fig, axes = plt.subplots(1, 3, figsize=(18,6))
for ax, (title, text) in zip(axes, texts.items()):
  wc = WordCloud(max_words=30).generate(text)
  ax.imshow(wc); ax.set_title(title)
  ax.axis('off')
\end{lstlisting}

\subsection{Slope Chart \& Sankey}
\textbf{Slope Chart}: show changes over time / rankings.\\
\zhint{斜率图：展示两个时间点之间的变化}\\
\textbf{Sankey Chart}: visualize flow between nodes.\\
\zhint{桑基图：展示数据流向（如能源流、供应链）}
\begin{lstlisting}
import plotly.graph_objects as go
fig = go.Figure(data=[go.Sankey(
  node=dict(pad=15, thickness=20,
    label=["SrcA","SrcB","ProcX","Out1"]),
  link=dict(source=[0,0,1],target=[2,3,2],
            value=[8,4,2]))])
fig.show()
\end{lstlisting}

% ==================== SECTION 6 ====================
\section{Image Parsing \& Model Architecture}

\subsection{Object Detection (Mask R-CNN)}
\zhint{目标检测：识别图片中的物体位置和类别}
\begin{lstlisting}
import torchvision
model = torchvision.models.detection\
  .maskrcnn_resnet50_fpn(weights=...)
model.eval()
with torch.no_grad():
  predictions = model([image_tensor])
# predictions: boxes, scores, labels, masks
\end{lstlisting}

\subsection{torchinfo -- Visualize Model Layers}
\zhint{查看神经网络每层的输出形状和参数数量}
\begin{lstlisting}
from torchinfo import summary
summary(model, input_size=(64, 28*28))
# Shows: layer type, output shape, param count
# Total params, trainable params, size
\end{lstlisting}

% ==================== SECTION 7 ====================
\section{GeoPandas \& Map Visualization}
\zhint{GeoPandas：在pandas基础上加地理空间数据支持}

\subsection{Loading Geospatial Data}
\begin{lstlisting}
import geopandas as gpd
# Load from file
gdf = gpd.read_file("example.shp")   # shapefile
gdf = gpd.read_file("example.geojson") # GeoJSON
gdf = gpd.read_file("example.gpkg")  # GeoPackage
# Load from URL
url = "https://naciscdn.org/naturalearth/..."
world = gpd.read_file(url)
\end{lstlisting}
\zhint{Shapefile=.shp+.shx+.dbf; GeoJSON=JSON格式; GPKG=SQLite}

\subsection{Geometry Attributes \& Methods}
\textbf{Attributes}: \texttt{area}, \texttt{bounds}, \texttt{total\_bounds}, \texttt{geom\_type}, \texttt{is\_valid}\\
\textbf{Methods}: \texttt{distance()}, \texttt{centroid}, \texttt{to\_crs()}, \texttt{plot()}

\subsection{Choropleth Map}
\zhint{分级统计图：用颜色深浅表示数值大小}
\begin{lstlisting}
# Simple world map
world.plot(figsize=(10,6))
# Choropleth with color by column
ax = world["geometry"].boundary.plot(
  figsize=(20,16))
merged_df.plot(column="Discipline", ax=ax,
  cmap='OrRd', legend=True,
  legend_kwds={"label":"Participation",
    "orientation":"horizontal"})
ax.set_title("Title"); plt.show()
\end{lstlisting}

\subsection{Layered Map}
\zhint{叠加多层地图：边界+区域+点+线}
\begin{lstlisting}
# Base layer: county boundaries
ax = counties.plot(figsize=(10,10),
  color='none', edgecolor='gainsboro')
# Add layers
wild_lands.plot(color='lightgreen', ax=ax)
campsites.plot(color='maroon',markersize=2,
               ax=ax)
trails.plot(color='black',markersize=1,ax=ax)
plt.show()
\end{lstlisting}

% ==================== SECTION 8 ====================
\section{Visualization Best Practices}
\zhint{可视化7条建议}
\begin{enumerate}[leftmargin=10pt,itemsep=0pt,topsep=0pt]
\item \textbf{Label directly} on the chart \zhint{直接在图上标注}
\item \textbf{Repeat units} \zhint{重复标注单位}
\item \textbf{Move axis ticks} where needed \zhint{刻度放在需要的地方}
\item \textbf{Don't center-align} text \zhint{不要居中对齐文字}
\item \textbf{Don't rotate} text \zhint{不要旋转文字}
\item \textbf{Use text outline} for readability \zhint{用文字描边提高可读性}
\item \textbf{Choose suitable number format} \zhint{选合适的数字格式}
\end{enumerate}

\keybox{
\textbf{Tools}: Matplotlib, Seaborn, Plotly, WordCloud, GeoPandas, torchinfo.\\
\textbf{Business}: Excel, Google Data Studio, Tableau, Power BI.
}

\end{multicols}
\end{CJK}
\end{document}
